\section{Related Work}
\label{sec:related}

The scheme sits at the intersection of three bodies of work: carbon-aware workload shifting, the cloud--edge continuum, and serverless cold-start mitigation. We position the scheme against the closest neighbours in each.

\paragraph{Carbon-aware workload shifting.}
The carbon-aware shifting literature is recent and expanding. Asadov et al.\ review the field and propose a spatio-temporal shifting algorithm, but the formulation uses the average signal and treats the site set as static~\cite{asadov2025carbonaware}. Patel et al.\ argue for an agile pathway to carbon-aware clouds on the marginal signal in principle but stop short of an online scheduler that consumes it~\cite{patel2024agile}. Radhakrishnan uses deep reinforcement learning to place Kubernetes workloads on a marginal signal but treats the workload set as coarse-grained batch jobs rather than per-invocation FaaS~\cite{radhakrishnan2024carbonaware}. Caribou shifts serverless applications geospatially between regions using marginal intensity but does not model temporal deferral across a renewable-surplus window~\cite{gsteiger2024caribou}. Rodrigues et al.\ apply temporal deferral to inter-datacentre data transfers, which lack the per-invocation deadline and cold-start structure of FaaS~\cite{rodrigues2025carbonaware}. Ozkan and Ozkan present a federated carbon-intelligence framework for AI workloads~\cite{ozkan2025federated}, and Shingne designs an iterative predictive carbon-aware scheduler for cloud--fog--edge~\cite{shingne2025design}; both emphasise prediction but neither gives a competitive-ratio guarantee against an offline optimum that knows the marginal trajectory. Lechowicz et al.\ give learning-augmented competitive algorithms for spatiotemporal online allocation with deadlines~\cite{lechowicz2024learningaugmented}, the closest theoretical neighbour, but their setting does not couple the marginal carbon signal to the per-invocation reuse decision that drives embodied amortisation. \textsc{CarbonShift} differs by consuming the marginal signal online, by exposing the embodied term through the reuse decision, and by giving an online guarantee against the marginal-aware offline optimum.

\paragraph{Carbon signals and metrics.}
The carbon-signal literature motivates the marginal choice. Sukprasert et al.\ measure the divergence between average and marginal signals and show that it changes the ranking of deferral targets~\cite{sukprasert2024implications}; Maji et al.\ extend this to the location/market-estimation dichotomy and show that the optimisation efficacy depends on the method~\cite{maji2024green}; Bashir et al.\ argue the sunk-carbon critique that operational-only accounting misallocates deferrals~\cite{bashir2024sunk}. These works establish the diagnosis; \textsc{CarbonShift} is the scheduler that acts on it. The flow-based grid carbon-accounting framework of Chen et al.\ gives the marginal-intensity estimation that the controller's signal rests on~\cite{chen2024towards}, and the renewable-forecasting literature provides the surplus-window forecast~\cite{singh2025riskaware}.

\paragraph{Cloud--edge continuum.}
Continuum modelling work~\cite{patel2024modeling,aldulaimy2024computing} formalises the computing continuum and its heterogeneity, and service-placement work on the continuum~\cite{rac2024costaware,soumplis2024performance,bellavista2024exploiting} formalises the cost of being in the wrong tier. Russo Russo et al.\ give QoS-aware offloading policies for serverless functions across the cloud-to-edge continuum~\cite{russo2024qosaware}; Calavaro et al.\ survey serverless functions in the compute continuum~\cite{calavaro2025beyond}. The continuum background supplies the eligible-site abstraction on which \textsc{CarbonShift}'s placements are defined; the carbon signal is the new dimension added to that abstraction.

\paragraph{Serverless cold starts and scheduling.}
The cold-start literature supplies the reuse decision that drives embodied amortisation. RainbowCake caches layers~\cite{yu2024rainbowcake}; CodeCrunch predicts invocations and warmups selectively with a cost model on the warmup~\cite{roy2024codecrunch}; cross-edge caching places functions probabilistically across edges~\cite{chen2024crossedge}; deep reinforcement learning orchestrates FaaS across the continuum~\cite{khansari2024deep}; Alatawi optimises multitenant resource allocation under reinforcement learning~\cite{alatawi2025optimizing}; Zhao et al.\ study the OS scheduler choice in FaaS~\cite{zhao2024serverless}; Hong et al.\ study the autoscaling/QoS interplay~\cite{hong2024analysis}; Chen et al.\ survey cold-start optimisation~\cite{chen2026cold}; Alelyani et al.\ give a fault-tolerant microservice scheduling strategy~\cite{alelyani2024optimizing}. None of these threads charges the cold-start decision with a carbon cost; the warmup that mitigates latency also amortises embodied carbon, and \textsc{CarbonShift} is the first scheduler to make that coupling explicit.

\paragraph{Hardware lifecycle and green ICT.}
The lifecycle literature supplies the embodied-cost terms. Schneider et al.\ give a cradle-to-grave lifecycle emissions model for AI hardware with generational trends~\cite{schneider2025lifecycle}; Ma and Zhou give a quantitative lifecycle cost and environmental-impact model for computing infrastructure~\cite{ma2024comprehensive}; the green-AI thematic analysis of Raman et al.\ situates sustainability in the broader AI research agenda~\cite{raman2024green}; Fernando and L\u{a}z\u{a}roiu survey energy-efficient IIoT in green 6G~\cite{raman2024energy}; Gill et al.\ give the modern-computing vision and challenges~\cite{gill2024modern}. The embodied terms in Eq.~\eqref{eq:cost} are drawn from these models.

\paragraph{Positioning.}
\textsc{CarbonShift} is the first per-invocation, marginal-carbon-aware, embodied-amortising, SLA-safe, non-gameable serverless deferral algorithm with an online competitive-ratio guarantee. The contributions are the coupling of the marginal signal to the per-invocation decision, the exposure of the embodied term through the reuse decision, and the three theorems that the scheme is provably competitive, SLA-safe, and non-gameable.
